\chapter{变系数扩散的守恒离散与对称性}
\label{app:variable-coeff}

\begin{keybox}
\textbf{目的：}从共享 face flux 出发重看变系数扩散，明确同一个守恒算子怎样同时具有 flux-difference FD 与 cell-centered FV 两种解释，并说明 face coefficient、harmonic average、守恒性与矩阵对称性之间的关系。
\end{keybox}


考虑一维
\begin{equation}
-\frac{d}{dx}\left(\beta(x)\frac{d\phi}{dx}\right)=f.
\end{equation}
在规则 cell-centered 网格上，下面的构造可以沿两条路线阅读。FV 路线把方程对控制体积分，并以 face flux 闭合；守恒型 FD 路线直接用相邻半点通量的差分近似散度。只要两条路线使用同一组 face flux，它们会生成同一局部代数行。这里把重点放在这个共享守恒算子的结构，以及两种推导如何到达同一代数结果。

对 cell center $i$，先在 face 定义通量
\begin{equation}
F_{i+1/2}
=-\beta_{i+1/2}\frac{\phi_{i+1}-\phi_i}{h}.
\end{equation}
FV 中对控制体收支求和、再除以 cell measure；守恒型 FD 中直接对两个半点通量作差。由于这里定义
\begin{equation}
F=-\beta\phi_x,
\end{equation}
原方程 $-(\beta\phi_x)_x=f$ 就是 $F_x=f$，因此两条路线都得到
\begin{equation}
\frac{F_{i+1/2}-F_{i-1/2}}{h}=f_i.
\label{eq:appb-flux-balance}
\end{equation}
代入 face flux 后整理得到
\begin{equation}
(A\phi)_i
=\frac{1}{h^2}
\left[
(\beta_{i-1/2}+\beta_{i+1/2})\phi_i
-\beta_{i-1/2}\phi_{i-1}
-\beta_{i+1/2}\phi_{i+1}
\right].
\end{equation}
若所有 face coefficient $\beta_{i+1/2}>0$，并且相邻 cells 对同一 face 使用同一个 coefficient，则该守恒离散矩阵保持对称，并在适当边界条件下为正定或半正定。这个结论来自共享耦合系数的代数结构，不依赖我们把该算子称为 FD 还是 FV。

当 $\beta$ 在相邻两个 cell 内近似分段常数，并且 face 位于两个 cell center 的中点时，可以把 center-to-center 路径看成两个长度均为 $h/2$ 的串联扩散阻抗。要求通过两段的稳态 flux 相同，可得到等效 face coefficient
\begin{equation}
\beta_{i+1/2}
=
\frac{2\beta_i\beta_{i+1}}
{\beta_i+\beta_{i+1}},
\end{equation}
即 harmonic average。

这里要注意适用边界：这个公式隐含了等距的两个半 cell。若界面位置偏离 face、网格非均匀或 coefficient reconstruction 更复杂，正确的 transmissibility 应使用实际几何距离和材料分段重新推导，而不是机械套用 $2\beta_i\beta_{i+1}/(\beta_i+\beta_{i+1})$。

对 GPU 而言，变系数 stencil 会增加 coefficient loads；对 AMG 而言，强跳变还会改变 strength graph 与 coarse-space 需求。

\section*{推导检查}
\begin{checkpointbox}
从 $F=-\beta\phi_x$ 出发重新推导式~\eqref{eq:appb-flux-balance}，检查通量差分前的符号。为什么相邻两个 cell 必须共享同一个 face coefficient 才容易保持矩阵对称？对位于两个 cell center 中点的材料界面，用“两个半 cell 串联阻抗”重新推导 harmonic average。
\end{checkpointbox}
